\clearpage
\section{Additional Work}
\label{sec:mindllm-additional-work}

\paragraph*{PRISM: Structured Text for Visual Reconstruction.}

\prism{} studies how the choice of intermediate representation affects reconstruction of viewed images from fMRI, and how object structure can improve the resulting images. It compares representations from language, vision, and vision-language models using centered kernel alignment, canonical correlation analysis, and the generalization gap of learned mappings from fMRI. Among the evaluated representations, caption-derived language-model embeddings show stronger alignment with the neural recordings, motivating a structured text interface between brain signals and image generation. A vision-language model generates training descriptions that separate objects, attributes, spatial locations, and relationships. An attribute and relationship search procedure selects descriptions that support reconstruction while meeting a neural-alignment constraint. An fMRI encoder and language model learn to predict these descriptions, which guide an object-centric diffusion process that composes separately conditioned object representations with global scene context. Experiments on NSD, BOLD5000, and GOD assess perceptual, structural, and semantic reconstruction quality, reporting reductions of up to 6\% in perceptual loss relative to prior methods. Ablations support both object-centric generation and attribute/relationship search. This work complements \mindllm{}'s semantic decoding by examining which representation spaces support neural alignment and how structured semantic content can guide visual reconstruction.

\paragraph*{D-Edit: Object-Level Control for Image Editing.}

Disentangled-Edit (\dedit{}) addresses object-level image editing: modifying a selected object while preserving the surrounding scene and maintaining consistency with text or reference-image guidance. The framework segments an image into editable items, including background regions, and assigns each item a distinct prompt composed of learned tokens. Grouped cross-attention restricts each item's interaction to its associated prompt, separating the control paths that would otherwise couple multiple objects. A two-stage adaptation first optimizes the new token embeddings and then fine-tunes diffusion-model cross-attention parameters to reconstruct the input. This procedure uses a single image, or a target and reference pair, without requiring captions of the original images. Users subsequently edit prompts, masks, or their associations to replace objects, refine appearance, move or reshape items, and remove unwanted content. Implementations with Stable Diffusion and SDXL are evaluated on text-guided and image-guided benchmarks containing 1,000 item-prompt combinations and 400 item-reference pairs. Quantitative comparisons show improved source-image preservation and guidance consistency over the evaluated baselines. Qualitative examples examine mask editing, while a user study assesses object removal. \dedit{} complements \mindllm{}'s interest in semantic interfaces by showing how generative controls can be localized to individual objects.
